\section{Experiments}
\label{sec:experiments}

The experiment asks two questions: do contemporary imitation policies
--- ACT~\citep{zhao2023aloha}, Diffusion
Policy~\citep{chi2023diffusion}, and the vision-language-action model
$\pi_{0.5}$~\citep{pi05} --- learn the two tasks from oracle
demonstrations collected through the teleoperation pipeline; and does
the simple/full mode split localise failures correctly, with simple
mode catching pipeline bugs and full mode measuring generalisation.
This section fixes the protocol; the results subsection is populated
from the long-run report as the living document grows.

\paragraph{Protocol.}
The campaign runs per mode (simple, then full), per task, per policy:
\begin{enumerate}
  \item \textbf{Oracle gate.} Both oracles are dry-run on thirty
    scenarios per task. The experiment aborts if the
    teleoperation-pipeline oracle drops below \SI{95}{\percent} on the
    single task or \SI{70}{\percent} on the relay. The relay bar sits
    below the oracle's measured per-attempt band (roughly
    \SIrange{75}{86}{\percent}) so a thirty-episode sample's noise
    cannot false-abort, yet far above the failure regime of a broken
    oracle.
  \item \textbf{Collection.} Simple mode collects repeated demonstrations
    of a single fixed scenario, which must therefore be one the oracle
    completes at a high per-attempt rate; the relay's per-scenario
    reliability varies, so its fixed scenario is selected by a short seed
    search rather than left at an arbitrary default. Full mode instead
    spreads collection across many training seeds and retries a failed
    seed a bounded number of times, so an occasional hard scenario is
    tolerated. Failed episodes are discarded and only verified successes
    saved.
  \item \textbf{Training.} ACT and Diffusion Policy train from scratch;
    $\pi_{0.5}$ is adapted from its published base with low-rank
    adapters, full finetuning being out of memory reach
    (Section~\ref{sec:diff-tooling}). All training uses the
    LeRobot framework~\citep{cadene2024lerobot}.
  \item \textbf{Checkpoint selection.} Every saved checkpoint rolls out
    on the validation seeds; the best checkpoint per cell (ties to the
    later step) advances.
  \item \textbf{Evaluation.} The selected checkpoint runs once per
    evaluation seed --- thirty held-out scenarios in full mode, thirty
    trials of the fixed scenario in simple mode --- with a video saved
    per episode. Success uses the collection criterion unchanged.
\end{enumerate}

\paragraph{What simple mode measures.}
Simple mode is a weaker instrument than its success rates suggest, and
the collected data say so directly: across every demonstration in a
simple-mode dataset the first recorded state is identical, its spread
zero on every channel. The mode also draws its validation and
evaluation scenarios from the seed it trained on, so a simple-mode
figure reports how completely a policy has fitted one scene and carries
no evidence about generalisation. We therefore read simple mode only as
the pipeline check it was introduced to be, and rely on full mode --- and
on the split-reach relay, whose cube spawn varies by construction --- for
any claim about generalisation.

\paragraph{Results.}
\textit{TODO: populate from the long-run report (success-rate table per
mode\,$\times$\,task\,$\times$\,policy with placement-error statistics,
per-seed success maps, validation curves, and head-to-head episode
renders) once the full campaign completes on the training machine. The
plumbing of every cell of this protocol --- collection, training,
checkpoint selection, evaluation, reporting --- has been verified end to
end at smoke scale.}
